% !TeX root = ../main.tex
\section{选题背景与意义}

\subsection{选题背景}

青年是整个社会力量中最积极、最有生气的力量\cite{Mao1991}，大学生的价值取向关乎个体成长方向，更决定着未来整个社会的价值追求与精神面貌\cite{Xi2014}。步入高等教育阶段的青年正处于价值观确立与定型的关键期，伴随生活空间拓展与社会接触面延伸，他们在学业规划、职业选择、日常消费、同伴互动与公共参与中不断深化对个人与社会关系的认识，逐步确立是非评判与行为抉择的基本尺度\cite{CPCC2017_Plan}。引导大学生扣好人生的“第一粒扣子”、自觉树立和践行社会主义核心价值观，构成了高校落实立德树人根本任务的核心环节\cite{CPCC2017_Ipe}。围绕培养德智体美劳全面发展的社会主义建设者和接班人这一总体目标\cite{Xi2018_Speech}，党的二十届三中全会对完善立德树人机制、健全培育和践行社会主义核心价值观制度机制作出系统部署\cite{CPCC2024_Jueding}，要求新时代高校价值观教育更加敏锐地切中青年的现实生存境遇。按照马克思主义唯物史观的根本观点，青年的现实生存境遇是由其所处的具体社会关系织就的；人的本质在其现实性上是一切社会关系的总和，受一定生产方式制约的现实交往形式直接规定着个体的精神过程与价值观念的生成\cite{MarxEngels2009}。大学生对主流价值的认知、检验与内化，始终依托于日常人际交往这一传递社会规范、交换生活经验与校准自我定位的现实中介。随着数字媒介深度融入校园生活以及社会分工持续细化，当代大学生的日常互动正突破传统熟人社会的差序格局，围绕细分兴趣、具体任务与阶段性诉求，同不同专业、院校及社会背景的个体建立起跨圈层联系。在家庭、行政班级与宿舍等稳定的强关系共同体之外，跨校学习伙伴、临时项目成员、趣缘社群熟人及匿名社区网友等“轻社交”与“搭子社交”日趋普遍。这类交往在情感卷入、背景了解与持续投入上相对有限，呈现出鲜明的弱连接特征。广泛分布于线上线下场景中的弱连接网络，承载着多元异质的生活经验与评价标准，成为审视当代大学生价值观生成与变迁的重要现实视域。

跨圈层弱连接之所以能够深度介入大学生的价值判断，源于异质性信息输入与低约束互动规则的叠加效应。不同生活背景与发展轨迹的交往对象充当着跨越群体边界的“信息桥”，在临时协作与垂直社群交流中将熟人圈层之外的生存经验引入大学生的认知视野，为其跳出封闭思维定势、培育现代契约意识与自主边界意识提供了客观条件。一旦这种缺乏深层情感连带与历史脉络支撑的浅表交往受到商业平台流量算法、消费主义景观及同辈竞争压力的渗透，跨圈层互动便极易滑向单向度的符号展演与功利算计，引发大学生在学业职业、生活消费、人际交往与社会公共领域的实然价值取向同唯物史观价值立场、社会主义核心价值观及当代育人要求之间的深层矛盾。社交平台高频推送的极端暴富样本与精细包装的优渥生活景观，在缺乏真实劳动过程参照的弱连接圈层中触发剧烈的向上社会比较，诱使部分学生在学业规划与择业抉择中摇摆于功利速成与消极退缩之间，或将自我价值寄托于商品符号的即时占有，陷入工具理性膨胀与消费异化的困境。日常同伴互动中按需匹配、随时断联的契约设计，将防卫性的投入产出核算引入人际伦理领域，把完整的社会关系拆解为单一诉求的短期互利交换，侵蚀了集体生活所需的道义承诺与信任根基。匿名网络公共讨论中剥离了具体历史语境与制度背景的情绪化表态借由弱连接网络快速传染，常将复杂的社会矛盾简化为二元对立标签，削弱青年作为历史创造者的主体担当以及理性、公正、法治的现代公民素养。

大学生日常交往形态向弱连接的延伸及其伴随的多维价值矛盾，要求高校思想政治教育深度切入学生的现实关系网络与微观生活世界。《关于新时代加强和改进思想政治工作的意见》明确提出把思想政治工作贯通于教育教学与日常管理全过程，贴近青年实际做好人文关怀与心理疏导\cite{CPCC2021_Yijian}。这一制度要求植根于马克思主义关于交往实践与道德养成的科学论断；列宁曾深刻指出，青年唯有将学习、教育和训练中的每一步骤同广泛的社会实践及共同体行动紧密结合，才能摆脱狭隘眼界的束缚，养成符合历史发展要求的道德品质\cite{Lenin2017}。立足数字化时代的育人环境，《教育强国建设规划纲要（2024—2035年）》进一步明确要坚持显性教育与隐性教育相统一，拓展实践育人与网络育人空间，统筹推进价值引领、实践体验与环境营造\cite{CPCC2025_Gangyao}。当前高校思想政治工作多依托院系、行政班级与宿舍等固定建制开展熟人式管理，面对跨越校园围墙、依即时诉求聚散的弱连接社交网络，往往存在感知滞后与覆盖盲区。本研究立足唯物史观关于社会关系与交往实践的基本原理，考察弱连接交往重塑大学生日常价值判断的作用机理与深层成因，探索切合青年交往特征的协同育人路径，为高校在流动社交环境中实现主流价值的有效引领提供学理支撑。

\subsection{研究意义}

本研究将大学生价值观的生成与变迁置于日常交往结构转型的具体情境中加以考察，有助于深化对思想政治教育过程中社会关系中介及其作用机理的学理认识。马克思主义关于现实的人及其社会交往的学说表明，个体价值观念的内化与外化始终受制于特定的社会关系结构与日常交往实践。既有学术探讨多将社会学的弱关系网络分析与思想政治教育学的价值观培育置于平行轨道，前者侧重考察跨圈层连接在传递异质信息与促成职业流动中的工具性效用，往往将行动者的价值观念视为外生给定的前提；后者偏重探讨主流意识形态的宏观阐释与话语传播，较少从历史唯物主义视角拆解数字弱连接改写青年具体价值判断的微观作用链条。本研究贯通社会网络结构分析与政治经济学批判，确立审视大学生日常价值变迁的规范标尺，深入剖析媒介算法规则与低约束互动契约转化为个体价值认知偏移的内在过程，拓展了弱连接理论在青年价值观研究领域的解释边界，为思想政治教育学把握日常社会关系变迁与主流价值内化的辩证关联提供了理论支点。

剖析弱连接交往引发的价值分歧及其生成逻辑，有助于高校准确把握大学生在跨圈层互动中的认知困惑，提升新时代思想政治教育的针对性。习近平总书记强调，做好高校思想政治工作要因事而化、因时而进、因势而新，密切关注青年在日常学习生活与社会交往中的现实关切\cite{Xi2022}。当代大学生关于择业标准、消费符号、人际边界与公共参与的观念波动，大量孕育于算法推送、匿名讨论与“搭子”协作等日常社交经历之中。透过表层社交行为识别学生所采纳的隐性评价标准及其背后的商业资本诱导、同辈竞争压力与信息视野局限，能够为思想政治理论课教学改革与日常学生管理提供鲜活的经验支撑。思想政治理论课作为落实立德树人根本任务的关键课程\cite{Xi2020}，需要将青年在现实交往中遭遇的价值困惑引入课堂教学，在回应学生深层关切的过程中落实《关于深化新时代学校思想政治理论课改革创新的若干意见》的部署要求\cite{CPCC2019_Yijian}。围绕成因机制所构建的引导方案，可直接服务于高校党团组织建设、跨学科科研协作、社会实践育人及校园网络空间治理的优化设计，推动思想政治教育实现叙事转换与文化浸润，将主流价值融入青年易于理解与参与的日常生活叙事\cite{Shen2021_narrative}，引导大学生运用唯物史观与政治经济学方法辨析碎片化信息，在将个人奋斗融入国家现代化建设的实践中实现自由全面发展。

